%%%% CAR-bench Competition Paper — IJCAI-ECAI 2026
%%%% Track: CAR-bench Competition (non-archival, 4 pages + refs)
%%%% Deadline: 26 July 2026 (AoE)

\documentclass{article}
\pdfpagewidth=8.5in
\pdfpageheight=11in

\usepackage{ijcai26}
\usepackage{times}
\usepackage{soul}
\usepackage{url}
\usepackage[hidelinks]{hyperref}
\usepackage[utf8]{inputenc}
\usepackage[small]{caption}
\usepackage{graphicx}
\usepackage{amsmath}
\usepackage{booktabs}
% Comment out for camera-ready:
\usepackage[switch]{lineno}
\linenumbers

\urlstyle{same}

\pdfinfo{
/TemplateVersion (IJCAI.2026.0)
}

\title{Deterministic Gates for LLM Agents:\\
       Improving Pass\textsuperscript{3} Consistency on CAR-bench}

% TODO: fill in author/affiliation before submission
\author{
    Author Name
    \affiliations
    Affiliation
    \emails
    email@example.com
}

\begin{document}

\maketitle

\begin{abstract}
% TODO: write after results are final (~200 words)
% Cover: thesis (deterministic shell raises Pass^3), method (7-stage pipeline),
% key result (Hallucination Pass^3 X→Y, Disambiguation X→Y, Base held),
% conclusion (deterministic pre/post-flight checks reduce variance without
% sacrificing capability).
\end{abstract}

% -----------------------------------------------------------------------
\section{Introduction}
% ~0.5 col — motivation + thesis
%
% Context: CAR-bench measures LLM agent consistency via Pass^3 (task passes
% all three independent runs). Stochastic LLM decisions destroy Pass^3 even
% when the agent "knows" the right answer. This is the gap we exploit.
%
% Thesis: A deterministic decision shell — capability checking, policy
% enforcement, fabrication guard — reduces variance and raises Pass^3,
% especially on Hallucination and Disambiguation, without hurting Base.
%
% Table: Baseline numbers (Tab.~\ref{tab:results}) show opportunity:
% Hallucination=0.48, Disambiguation=0.46 for best-known agent (Claude Opus).

\section{Method: Deterministic Shell Architecture}
% ~1.5 col — the 7-stage pipeline and why each stage matters

\subsection{Ledger and State Machine}
% Stages 1--2: provenance ledger + INTAKE→CAPABILITY\_CHECK→(CLARIFY|PLAN)
% →POLICY\_CHECK→EXECUTE→VERIFY→RESPOND with bounded loop, idempotent call IDs

\subsection{Capability Matcher}
% Stage 3: deterministic 3-way check at every planning step (not just intake)
% Three removal types: tool, parameter, result-field

\subsection{Policy Compiler}
% Stage 4: 19 policies as pre-flight predicates; structural impossibility
% of violations rather than probabilistic avoidance

\subsection{Fabrication Guard}
% Stage 5: every claim in the response needs provenance in the ledger
% (tool observation, user statement, policy default)

\subsection{Disambiguation Engine}
% Stage 6: internal resolution first (preferences, policies, context);
% user query only when genuinely unresolvable + action is state-changing

\subsection{Auditor}
% Stage 7: targeted self-check at two expensive checkpoints only

% -----------------------------------------------------------------------
\section{Experiments}

\subsection{Setup}
% Agent model: anthropic/claude-sonnet-4-6 (dev), claude-opus-4-7 (eval)
% User-Sim + Judge: gemini-2.5-flash
% Splits: base (100 train / 100 public-val), hallucination (98/100),
%         disambiguation (56/69)
% Metric: Pass^3 (3 independent trials, task counts only if all three pass)

\subsection{Calibration Run (10 July)}
% TODO: fill in after calibration run
% Report Pass^3 on public-validation set, compare to baseline

% -----------------------------------------------------------------------
\section{Results}

\begin{table}[t]
\centering
\begin{tabular}{lrrrr}
\toprule
Agent & Overall & Base & Hall. & Disamb. \\
\midrule
Claude Opus 4.6 (vanilla) & 0.58 & 0.80 & 0.48 & 0.46 \\
GPT-5 thinking            & 0.54 & 0.66 & 0.60 & 0.36 \\
\midrule
Ours (TODO)               & --   & --   & --   & --   \\
\bottomrule
\end{tabular}
\caption{Pass\textsuperscript{3} on public-validation set.
         Baseline numbers are publicly reported.
         Our result will be filled in after the final evaluation (19 July).}
\label{tab:results}
\end{table}

% TODO: ablation table (each stage's contribution to Pass^3)
% TODO: Pass^1 / Pass^3 / Pass@3 gap analysis (shows where variance was)

% -----------------------------------------------------------------------
\section{Related Work}

% CAR-bench~\cite{kirmayr2026carbench}: the benchmark we evaluate on (required cite)
% tau-bench~\cite{yao2024taubench}: tool-agent-user interaction benchmark
% Deterministic agent shells / constrained LLM decoding (related approaches)

% -----------------------------------------------------------------------
\section{Conclusion}
% 3--4 sentences: thesis confirmed/refuted, key numbers, open limitations

\section*{Acknowledgements}
% optional — add if applicable

\bibliographystyle{named}
\bibliography{references}

\end{document}
